% chapter_lqr.tex  —  LQR Control for Gust Load Alleviation
% Standalone compilable chapter. Compile with: pdflatex chapter_lqr.tex
\documentclass[11pt,a4paper]{article}
\usepackage{amsmath, amssymb}
\usepackage{bm}           % \bm{} for bold math
\usepackage{graphicx}
\usepackage{subcaption}
\usepackage{booktabs}
\usepackage{hyperref}
\usepackage[margin=2.5cm]{geometry}

\newcommand{\vxi}{\bm{\xi}}
\newcommand{\vu}{\bm{u}}
\newcommand{\vy}{\bm{y}}
\newcommand{\mA}{\mathbf{A}}
\newcommand{\mB}{\mathbf{B}}
\newcommand{\mC}{\mathbf{C}}
\newcommand{\mD}{\mathbf{D}}
\newcommand{\mK}{\mathbf{K}}
\newcommand{\mP}{\mathbf{P}}
\newcommand{\mQ}{\mathbf{Q}}
\newcommand{\mR}{\mathbf{R}}
\newcommand{\trim}{\mathrm{tr}}

\title{LQR Control for Gust Load Alleviation\\on a Neural Aeroelastic Model}
\author{}
\date{}

\begin{document}
\maketitle

% ================================================================
\section{Introduction}
\label{sec:intro}
% ================================================================

Gust load alleviation (GLA) aims to reduce the structural response of an
aeroelastic system to atmospheric disturbances through active flap deflection.
Model Predictive Control (MPC) provides excellent performance by rolling out
an explicit nonlinear preview over a receding horizon, but requires solving an
online optimisation at every timestep.
Linear Quadratic Regulator (LQR) offers an attractive alternative: the optimal
feedback gain is computed \emph{offline} once, at the expense of restricting
optimality to the neighbourhood of the trim point.

The plant considered here is the \emph{LDNet} model---a data-driven aeroelastic
emulator that couples a two-degree-of-freedom (2-DOF) structural wing with a
neural-network aerodynamic surrogate.
Because the aerodynamic model is a trained neural network, closed-form
Jacobians are unavailable; they are instead obtained through \emph{automatic
differentiation} (AD) of the full discrete-time propagator.
The resulting linearised system is then used to solve the Discrete Algebraic
Riccati Equation (DARE) and extract the optimal gain matrix $\mK$.

% ================================================================
\section{Aeroelastic Model and Trim Condition}
\label{sec:model}
% ================================================================

\subsection{Structural Dynamics}

The wing is modelled as a rigid section with two degrees of freedom: heave
displacement $h$ (positive downward) and pitch angle $\alpha$ (positive
nose-up).  Aerodynamic forces are transmitted to the structure via a trailing-edge
flap with deflection angle $\delta$.

The equations of motion in coupled form are
\begin{equation}
  \begin{bmatrix} M_{hh} & M_{h\alpha} \\ M_{h\alpha} & M_{\alpha\alpha} \end{bmatrix}
  \begin{bmatrix} \ddot{h} \\ \ddot{\alpha} \end{bmatrix}
  =
  \begin{bmatrix} -F_y - D_h\dot{h} - K_h h \\ M_z - D_\alpha\dot{\alpha} - K_\alpha\alpha \end{bmatrix},
  \label{eq:eom}
\end{equation}
where $F_y = q_\infty C_L$ and $M_z = q_\infty C_M$ are the aerodynamic lift
force and pitching moment, $q_\infty = \tfrac{1}{2}\rho U_\infty^2 S$ is the
dynamic pressure ($S = 0.05\,\text{m}^2$, $\rho = 1.225\,\text{kg/m}^3$), and
the inertial coupling arises from the off-centre flap mass:
\begin{equation}
  M_{hh} = 24.09\,\text{kg},\quad
  M_{\alpha\alpha} = 2.063\,\text{kg\,m}^2,\quad
  M_{h\alpha} = 0.625\,\text{kg\,m}.
\end{equation}
Structural stiffness and damping parameters are
$K_h = 4000\,\text{N/m}$, $K_\alpha = 700\,\text{N\,m/rad}$,
$D_h = 12\,\text{N\,s/m}$, $D_\alpha = 1.6\,\text{N\,m\,s/rad}$,
corresponding to modal damping ratios of approximately $2\%$.

\subsection{Neural Aerodynamic Surrogate (LDNet)}

The aerodynamic coefficients $C_L$ and $C_M$ are produced by a recurrent
neural network with a scalar latent state $z \in \mathbb{R}$:
\begin{align}
  z_{k+1} &= z_k + \frac{\Delta t}{\Delta t_{\rm ref}}\,f_{\rm dyn}(z_k, h_k, \dot{h}_k, \alpha_k, \dot{\alpha}_k, \delta_k, W_k;\,\theta_d),
  \label{eq:zdyn}\\
  \begin{bmatrix} C_L \\ C_M \end{bmatrix}_k &= f_{\rm rec}(z_k, h_k, \dot{h}_k, \alpha_k, \dot{\alpha}_k, \delta_k;\,\theta_r),
  \label{eq:yrec}
\end{align}
where $W_k$ is the gust velocity, $f_{\rm dyn}$ is a shallow MLP (7-7-1), and
$f_{\rm rec}$ is a deeper network (24-24-24-24-2).
Both networks are kept frozen during control design.

\subsection{Augmented State and Trim}

The full state is the concatenation of structural and latent variables:
\begin{equation}
  \vxi_k = \begin{bmatrix} h_k & \dot{h}_k & \alpha_k & \dot{\alpha}_k & z_k \end{bmatrix}^\top \in \mathbb{R}^5.
\end{equation}
The trim condition is defined at free-stream velocity $U_\infty = 75\,\text{m/s}$
with zero gust ($W = 0$) and zero structural displacement.
The latent state $z_{\rm tr}$ is obtained by iterating \eqref{eq:zdyn} to
steady state with $\delta = 0$; at trim $C_{L,\rm tr} \approx 0.031$ and
$C_{M,\rm tr} \approx 0$.
The trim point is therefore
\begin{equation}
  \vxi_{\rm tr} = \begin{bmatrix} 0 & 0 & 0 & 0 & z_{\rm tr} \end{bmatrix}^\top, \quad
  u_{\rm tr} = \delta_{\rm tr} = 0.
\end{equation}

% ================================================================
\section{Jacobian Linearisation via Automatic Differentiation}
\label{sec:jacobian}
% ================================================================

\subsection{Discrete-Time Propagator}

The combined structural integrator and aerodynamic network define a
discrete-time propagator
\begin{equation}
  \vxi_{k+1} = F(\vxi_k, u_k),
  \label{eq:propagator}
\end{equation}
where $F: \mathbb{R}^5 \times \mathbb{R} \to \mathbb{R}^5$ is evaluated in
two stages:
\begin{enumerate}
  \item \textbf{Aerodynamic step.}  Given $(\vxi_k, u_k)$, evaluate
        LDNet to obtain $z_{k+1}$, $C_L$, $C_M$ via
        \eqref{eq:zdyn}--\eqref{eq:yrec} with $W = 0$ (no gust at trim).
  \item \textbf{Structural step.}  Integrate \eqref{eq:eom} over one timestep
        $\Delta t = 0.01\,\text{s}$ using a fourth-order Runge--Kutta (RK4)
        scheme with the aerodynamic forcing fixed at the values from step~1:
        \begin{align}
          \dot{\bm{x}} &= \bm{g}(\bm{x}, F_y, M_z), \nonumber\\
          g_1 &= \dot{h},\quad
          g_2 = \frac{M_{\alpha\alpha}(-F_y - D_h\dot{h} - K_h h)
                      - M_{h\alpha}(M_z - D_\alpha\dot{\alpha} - K_\alpha\alpha)}
                     {M_{hh}M_{\alpha\alpha} - M_{h\alpha}^2},\nonumber\\
          g_3 &= \dot{\alpha},\quad
          g_4 = \frac{M_{hh}(M_z - D_\alpha\dot{\alpha} - K_\alpha\alpha)
                      - M_{h\alpha}(-F_y - D_h\dot{h} - K_h h)}
                     {M_{hh}M_{\alpha\alpha} - M_{h\alpha}^2},
          \label{eq:rhs}
        \end{align}
        where $\bm{x} = [h,\dot{h},\alpha,\dot{\alpha}]^\top$ and the
        determinant $\det \mathbf{M} = M_{hh}M_{\alpha\alpha} - M_{h\alpha}^2$.
\end{enumerate}
The output equation evaluated at the same step is
\begin{equation}
  \vy_k = \begin{bmatrix} C_L \\ C_M \end{bmatrix}_k = g(\vxi_k, u_k).
\end{equation}

\subsection{Four Jacobian Matrices}

Linearising \eqref{eq:propagator} and the output map around the trim point
$(\vxi_{\rm tr}, 0)$ yields the discrete-time LTI system
\begin{equation}
  \delta\vxi_{k+1} = \mA_d\,\delta\vxi_k + \mB_d\,\delta u_k, \qquad
  \delta\vy_k = \mC_y\,\delta\vxi_k + \mD_y\,\delta u_k,
  \label{eq:lti}
\end{equation}
where $\delta\vxi = \vxi - \vxi_{\rm tr}$ and $\delta u = u - u_{\rm tr}$, and
\begin{equation}
  \mA_d = \frac{\partial F}{\partial \vxi}\bigg|_{\rm tr} \in \mathbb{R}^{5\times5},\quad
  \mB_d = \frac{\partial F}{\partial u}\bigg|_{\rm tr} \in \mathbb{R}^{5\times1},\quad
  \mC_y = \frac{\partial g}{\partial \vxi}\bigg|_{\rm tr} \in \mathbb{R}^{2\times5},\quad
  \mD_y = \frac{\partial g}{\partial u}\bigg|_{\rm tr} \in \mathbb{R}^{2\times1}.
  \label{eq:jacobians}
\end{equation}

\subsection{Computation via \texttt{tf.GradientTape}}

Because $F$ passes through the LDNet weights, closed-form expressions for the
partial derivatives in \eqref{eq:jacobians} are not available.
Instead, all four matrices are computed in a single forward pass under a
TensorFlow persistent \texttt{GradientTape}:

\begin{enumerate}
  \item Declare $\vxi_{\rm tr}$ and $u_{\rm tr}$ as \texttt{tf.Variable}s so
        that the tape tracks operations on them.
  \item Evaluate $F(\vxi_{\rm tr}, u_{\rm tr})$ and $g(\vxi_{\rm tr}, u_{\rm tr})$
        inside the tape context.  Both the RK4 stages in \eqref{eq:rhs} and the
        neural-network layers of LDNet are executed as differentiable TensorFlow
        operations.
  \item After the forward pass, call \texttt{tape.jacobian} to obtain $\mA_d$,
        $\mB_d$, $\mC_y$, $\mD_y$ by reverse-mode automatic differentiation.
\end{enumerate}

This approach requires no analytical linearisation of the neural network and
correctly propagates gradients through the RK4 sub-stages as well as the
nonlinear activations of $f_{\rm dyn}$ and $f_{\rm rec}$.

The open-loop eigenvalues $\lambda_i(\mA_d)$ all satisfy $|\lambda_i| < 1$,
confirming that the trim equilibrium is asymptotically stable in the
linearised sense, albeit with low modal damping
($|\lambda_i| \approx 0.97$--$0.99$ for the aeroelastic modes).

% ================================================================
\section{LQR Formulation and Discrete Algebraic Riccati Equation}
\label{sec:lqr}
% ================================================================

\subsection{Cost Function with Output Penalty}

Given the LTI system \eqref{eq:lti}, we seek a state-feedback gain $\mK$
such that the infinite-horizon cost
\begin{equation}
  J = \sum_{k=0}^{\infty}
      \Bigl[\delta\vxi_k^\top \mQ_{\rm aug}\,\delta\vxi_k
           + \delta u_k^\top \mR_{\rm aug}\,\delta u_k\Bigr]
  \label{eq:cost}
\end{equation}
is minimised.  To penalise deviations of the aerodynamic outputs $[C_L, C_M]$
from trim without increasing the state dimension, the weight matrices are
augmented with the output Jacobians:
\begin{equation}
  \mQ_{\rm aug} = \mQ + \mC_y^\top \mQ_y \mC_y, \qquad
  \mR_{\rm aug} = \mR + \mD_y^\top \mQ_y \mD_y,
  \label{eq:augmented}
\end{equation}
where $\mQ \in \mathbb{R}^{5\times5}$ is a diagonal state cost,
$\mR \in \mathbb{R}^{1\times1}$ is the control cost, and
$\mQ_y \in \mathbb{R}^{2\times2}$ is the output penalty.
Equation~\eqref{eq:augmented} follows directly from substituting the output
equation into the cost: the cross terms
$\delta\vxi^\top\mC_y^\top\mQ_y\mC_y\delta\vxi$ appear in the augmented state
cost, and similarly for $\mD_y$.

The weight values used in this work are
\begin{equation}
  \mQ = \operatorname{diag}\!\left(\frac{1}{0.00428^2},\,0,\,\frac{5}{0.00823^2},\,0,\,0\right),
  \quad
  \mR = \left[\frac{1}{0.5^2}\right],
  \label{eq:weights}
\end{equation}
normalised by the expected peak amplitudes of heave and pitch under a
$60\,\text{m/s}$ gust.  No output penalty is applied ($\mQ_y = 0$) for the
baseline LQR presented here; the role of $\mQ_y$ is discussed in
Section~\ref{sec:sensitivity}.

\subsection{Discrete Algebraic Riccati Equation}

The solution to the infinite-horizon LQR problem is given by the unique
positive semi-definite solution $\mP$ of the \emph{Discrete Algebraic Riccati
Equation} (DARE):
\begin{equation}
  \mP = \mQ_{\rm aug}
        + \mA_d^\top \mP \mA_d
        - \mA_d^\top \mP \mB_d
          \bigl(\mR_{\rm aug} + \mB_d^\top \mP \mB_d\bigr)^{-1}
          \mB_d^\top \mP \mA_d,
  \label{eq:dare}
\end{equation}
which is solved numerically via \texttt{scipy.linalg.solve\_discrete\_are}
\cite{anderson1990optimal}.

\subsection{Optimal Gain and Control Law}

The optimal state-feedback gain is
\begin{equation}
  \mK = \bigl(\mR_{\rm aug} + \mB_d^\top \mP \mB_d\bigr)^{-1}
        \mB_d^\top \mP \mA_d
        \in \mathbb{R}^{1\times5},
  \label{eq:gain}
\end{equation}
and the control law applied at each timestep is
\begin{equation}
  \delta_t = -\mK\,(\hat{\vxi}_t - \vxi_{\rm tr}),\qquad
  \delta_t \in [-20^\circ,\,+20^\circ],
  \label{eq:law}
\end{equation}
where $\hat{\vxi}_t$ is the current state estimate.
The closed-loop matrix $\mA_d - \mB_d\mK$ has all eigenvalues strictly inside
the unit circle, confirming closed-loop stability.

The gain is computed once at design time; at runtime \eqref{eq:law} requires
only a single matrix-vector product and a scalar clip---negligible computational
cost compared to the online MPC optimisation.

\begin{thebibliography}{9}
\bibitem{anderson1990optimal}
B.~D.~O.~Anderson and J.~B.~Moore,
\textit{Optimal Control: Linear Quadratic Methods}.
Prentice Hall, 1990.
\end{thebibliography}

% ================================================================
\section{Simulation Results}
\label{sec:results}
% ================================================================

\subsection{Setup}

The LQR controller is evaluated against an open-loop baseline (no flap
actuation) under a single 1-cosine gust:
\begin{equation}
  W(t) = \frac{W_0}{2}\Bigl(1 - \cos\tfrac{2\pi t}{T_g}\Bigr),
  \quad 0 \le t \le T_g,
\end{equation}
with peak gust velocity $W_0 = 60\,\text{m/s}$ and duration $T_g = 1\,\text{s}$,
at free-stream speed $U_\infty = 75\,\text{m/s}$.
The structural state is initialised at trim; $z_0 = z_{\rm tr}$.
The state estimate $\hat{\vxi}_t$ is provided by a leaky kinematic integrator
on the measured heave acceleration and a pitch angle sensor (angle of attack
measurement).

\subsection{Structural Response}

Figure~\ref{fig:state} shows the time histories of heave displacement $h$,
heave velocity $\dot{h}$, pitch angle $\alpha$, and pitch rate $\dot{\alpha}$.
The LQR controller significantly damps the post-gust oscillations in both
channels.  Table~\ref{tab:metrics} summarises the amplitude reduction (peak-to-peak
amplitude over the window $[0,\,1.5]\,\text{s}$) for all key outputs.

\begin{figure}[htbp]
  \centering
  \includegraphics[width=\textwidth]{figures/fig1_state.png}
  \caption{Structural state time histories under a $60\,\text{m/s}$ 1-cosine gust.
           Blue: open-loop baseline. Cyan: LQR controlled response.}
  \label{fig:state}
\end{figure}

\subsection{Structural Accelerations}

Figure~\ref{fig:accels} shows heave and pitch accelerations, which are directly
related to structural loads.  The LQR achieves a marked reduction in peak pitch
acceleration, with a modest improvement in heave acceleration.

\begin{figure}[htbp]
  \centering
  \includegraphics[width=\textwidth]{figures/fig2_accels.png}
  \caption{Structural acceleration time histories.
           Blue: open-loop. Cyan: LQR.}
  \label{fig:accels}
\end{figure}

\subsection{Control Effort}

The flap deflection $\delta(t)$ and its rate $\dot{\delta}(t)$ are shown in
Figure~\ref{fig:control}.  The control input does not saturate (stays within
$\pm 20^\circ$), and the rate remains smooth, confirming that the linear
approximation is valid over the gust duration.

\begin{figure}[htbp]
  \centering
  \includegraphics[width=\textwidth]{figures/fig3_control.png}
  \caption{Flap deflection $\delta(t)$ and rate $\dot{\delta}(t)$ under LQR.
           Dotted lines: $\pm 20^\circ$ saturation limits.}
  \label{fig:control}
\end{figure}

\subsection{Aerodynamic Coefficients}

Figure~\ref{fig:aero} shows $C_L(t)$ and $C_M(t)$.  The initial $C_L$ peak at
$t \approx 0.3\,\text{s}$ is not eliminated by the LQR: because the controller
has no preview of the incoming gust, it reacts only \emph{after} the gust
excites the aerodynamics.  The subsequent oscillations, however, are noticeably
attenuated.

\begin{figure}[htbp]
  \centering
  \includegraphics[width=\textwidth]{figures/fig4_aero.png}
  \caption{Lift and pitching moment coefficient time histories.
           Blue: open-loop. Cyan: LQR.}
  \label{fig:aero}
\end{figure}

\subsection{Performance Summary}

\begin{table}[htbp]
  \centering
  \caption{Peak-to-peak amplitude in the gust window $[0,\,1.5]\,\text{s}$.
           Reduction relative to open-loop baseline.}
  \label{tab:metrics}
  \begin{tabular}{lcc}
    \toprule
    Quantity & Baseline & LQR  \\
    \midrule
    Heave displacement $|h|_{\rm pp}$ [mm]          & --- & --- \\
    Pitch angle $|\alpha|_{\rm pp}$ [$^\circ$]       & --- & --- \\
    Heave acceleration $|\ddot{h}|_{\rm pp}$ [m/s$^2$] & --- & --- \\
    Pitch acceleration $|\ddot{\alpha}|_{\rm pp}$ [$^\circ$/s$^2$] & --- & --- \\
    Lift coefficient $|C_L|_{\rm pp}$                & --- & --- \\
    Pitching moment $|C_M|_{\rm pp}$                 & --- & --- \\
    \bottomrule
  \end{tabular}
\end{table}

\noindent\textit{Fill numerical values from the terminal output of
\texttt{src/test\_mpc.py} before final submission.}

% ================================================================
\section{Sensitivity to Q/R Weights and Design Trade-offs}
\label{sec:sensitivity}
% ================================================================

The LQR design involves three tuning matrices: $\mQ$ (state cost), $\mR$
(control cost), and optionally $\mQ_y$ (output penalty).  Their influence on
closed-loop behaviour follows directly from the structure of the DARE solution.

\paragraph{Effect of $\mQ$ and $\mQ_y$.}
Increasing the heave weight $Q_h$ or output penalty $Q_y$ causes the Riccati
solution $\mP$ to grow, which in turn increases the magnitude of $\mK$ via
\eqref{eq:gain}.  A larger gain produces faster disturbance rejection but
requires higher flap deflections.  In the extreme case where $Q_h \gg R$, the
controller approaches a minimum-time regulator and $\delta$ may saturate the
$\pm 20^\circ$ limit, rendering the linear approximation invalid.

\paragraph{Effect of $\mR$.}
Increasing $\mR$ penalises control effort directly.  The gain $\mK$ decreases,
resulting in slower but smoother actuation.  For the present system the ratio
$Q_h / R$ is the dominant design parameter: at the baseline value
$Q_h/R \approx 54596 / 4 \approx 13649$, the controller achieves good
oscillation attenuation without saturating the actuator.

\paragraph{Output penalty $\mQ_y$.}
Setting $\mQ_y > 0$ shifts the cost from state excursion to aerodynamic
output deviation.  Because the output Jacobian $\mC_y$ couples all five states
to $C_L$ and $C_M$, the augmented weight $\mQ_{\rm aug}$ becomes full even when
$\mQ$ is diagonal.  This can improve $C_M$ attenuation at the cost of larger
heave excursions, trading structural load metrics against structural
displacement metrics.

\paragraph{Implicit sensitivity.}
The gain $\mK(\mQ,\mR)$ is an implicit function of the weight matrices through
the DARE solution.  No closed-form sensitivity expression exists; the
dependence must be evaluated numerically.  For extreme weight ratios
($Q/R > 10^6$ or $Q/R < 10^{-2}$), the DARE may become ill-conditioned,
and the solver may fail to converge.

% ================================================================
\section{Conclusions}
\label{sec:conclusions}
% ================================================================

This chapter has presented an LQR controller for gust load alleviation on the
LDNet aeroelastic model.  The key contributions are:

\begin{itemize}
  \item A linearisation procedure based on automatic differentiation that
        obtains exact Jacobians of a neural-network plant without manual
        derivation.
  \item An output-augmented LQR formulation that penalises aerodynamic
        coefficient deviations at trim using the output Jacobians
        $\mC_y$ and $\mD_y$.
  \item Demonstration that a single offline DARE solve provides a
        real-time-capable gain (one matrix-vector multiply per timestep)
        with meaningful structural load reduction under a $60\,\text{m/s}$
        1-cosine gust.
\end{itemize}

The primary limitation of the LQR approach is its reactive nature: without
preview information about the incoming gust, the initial aerodynamic peak
at $t \approx 0.3\,\text{s}$ cannot be avoided.  MPC, by rolling out an
explicit nonlinear horizon, achieves further reduction of this first peak
at significantly higher computational cost.  The LQR presented here represents
a practical baseline for real-time implementation where computational resources
are constrained.

\end{document}
